\documentclass[11pt,a4paper]{article}
\usepackage[margin=2.3cm]{geometry}
\usepackage{amsmath,amssymb,booktabs,microtype}
\usepackage[dvipsnames]{xcolor}
\usepackage[colorlinks=true,linkcolor=MidnightBlue,urlcolor=MidnightBlue]{hyperref}
\newcommand{\bq}{\mathbf{q}}
\newcommand{\br}{\mathbf{r}}
\title{Why the old large-cell $\Gamma$ scheme is not a reference for 2D GW:\\the spurious charge of fitted pair densities}
\author{Ritaj Tyagi, University of Regensburg}
\date{6 November 2026}
\begin{document}
\maketitle

\paragraph{Summary.}
The large-cell $\Gamma$-point GW implementations in CP2K (\texttt{tensor\_large\_cell\_Gamma},
\texttt{ri\_rs\_large\_cell\_Gamma}) build the screened interaction of a 2D supercell on a
k-mesh \texttt{KPOINTS\_W}~$n\,n\,1$ of the supercell's small Brillouin zone, with a second
``extrapolation'' mesh of $2n$ points and a Coulomb matrix $V(\bq)$ from a minimum-image
truncated real-space lattice sum. Band gaps from this scheme were used as the reference against
which the new $\Gamma$-only scheme (one real k-point plus an analytic head-and-wing correction)
was judged, and the two differ by $0.12$~eV on the largest cell we can afford. This note
documents that the old scheme carries the systematic error, that the error is the spurious net
charge of RI-fitted pair densities identified by J.~Wilhelm in cp2k PR~\#6447, and that the
$\Gamma$-only numbers are confirmed by his full-k-point code to about 15~meV. Numbers that came
out of the old scheme for 2D systems should therefore not be used as references.

\section{The observation}

Monolayer h-BN, $a = 2.502$~\AA, 15~\AA\ vacuum, PBE, aug-SZV-MOLOPT-GTH tier-1 orbital and RI
bases, 10 minimax points, \texttt{CUTOFF\_RADIUS\_RI}~5~\AA, identical in every code and run.
``Gap'' is the direct G$_0$W$_0$ gap at the k-point $P = (11/32, 11/32)$ of the primitive zone,
the point nearest to $K$ that an $8\times8$ supercell with a $2\times2$ DOS mesh samples.

\begin{center}
\begin{tabular}{llr}
\toprule
scheme & cell / mesh & gap (eV)\\
\midrule
old large-cell scheme, \texttt{KPOINTS\_W 4 4 1} + extrapolation & $8\times8$ supercell & 8.152\\
old large-cell scheme, \texttt{KPOINTS\_W 4 4 1} + extrapolation & $14\times14$ supercell & 6.582\\
old large-cell scheme, \texttt{KPOINTS\_W 2 2 1} + extrapolation & $14\times14$ supercell & 6.595\\
\midrule
$\Gamma$-only + head/wing correction (this work) & $8\times8$ supercell & 6.710\\
$\Gamma$-only + head/wing correction & $10\times10$ supercell & 6.705\\
$\Gamma$-only + head/wing correction & $14\times14$ supercell & 6.699\\
\midrule
full-k RI-RS GW, PR \#6447 (primitive cell) & $6\times6$ mesh & 6.590\\
full-k RI-RS GW, PR \#6447 & $8\times8$ mesh & 6.658\\
full-k RI-RS GW, PR \#6447 & $12\times12$ mesh & 6.681\\
full-k RI-RS GW, PR \#6447 & $16\times16$ mesh & 6.684\\
full-k RI-RS GW, PR \#6447 & $20\times20$ mesh & 6.687\\
full-k RI-RS GW, PR \#6447 & $24\times24$ mesh & 6.690\\
\bottomrule
\end{tabular}
\end{center}

Three things stand out. (i) The old scheme is not converged with the cell: $8.15$~eV at
$8\times8$ against $6.58$~eV at $14\times14$, and an earlier sweep with the tier-2 basis
(11.85, 9.74, 9.66, 9.43, 8.63, 7.54~eV for $n = 5, 6, 7, 8, 10, 12$) was still falling at
30~\AA. (ii) The $\Gamma$-only scheme is converged to 5~meV in the gap from $8\times8$ on.
(iii) The independent full-k-point code converges from below onto the $\Gamma$-only value;
state by state at the supercell $\Gamma$ point the absolute G$_0$W$_0$ energies agree to 5~meV
(valence) and 14~meV (conduction) with the $24\times24$ mesh, with identical $\Sigma^x$ and
$v^{xc}$ and $\Sigma^c$ within 10~meV. The old scheme sits $0.12$~eV below both at
$14\times14$ and is still drifting downwards, so the $0.12$~eV is its error, not the new
code's.

\section{The mechanism: a spurious charge in the fitted pair densities}

In the real-space RI (RI-RS) a product of two orbitals is interpolated from its values on the
atom-centred grid points $\br_\ell$ and fitted to the auxiliary basis through the fit
coefficients $Z_{\ell P}$,
\begin{equation}
  \psi_i(\br)\psi_a(\br) \approx \sum_{P}\Big[\sum_\ell Z_{\ell P}\,\psi_i(\br_\ell)\psi_a(\br_\ell)\Big]\varphi_P(\br)
  \equiv \sum_P d_{ia,P}\,\varphi_P(\br).
\end{equation}
The exact transition density is neutral, $\int\psi_i\psi_a = 0$ by orthogonality. The fitted one
is not: its charge $m_{ia} = \sum_P c_P\,d_{ia,P}$, with $c_P = \int\varphi_P$, is a small but
nonzero number set by the grid and the fit (the grid cannot represent orthogonality, and the
fit is regularised). This was found by J.~Wilhelm while converging his full-k-point RI-RS GW
(cp2k PR~\#6447, commit \texttt{a2a9da4c91}).

A spurious charge is harmless for a molecule and nearly harmless for a 3D crystal, but not in
two dimensions. On the $\bq$ mesh of a 2D calculation the Coulomb kernel in the charge channel
is $v(\bq) = 2\pi/(Aq)$, so the charge mode of $\chi$ enters the dielectric function as
\begin{equation}
  \varepsilon_{00}(\bq) - 1 \;\sim\; \frac{2\pi}{Aq}\sum_{ia}\frac{2\,|m_{ia}|^2\,\Delta_{ia}}{\Delta_{ia}^2+\omega^2}
  \;=\; \frac{C}{q},
\end{equation}
a $1/q$ divergence with a spurious prefactor. Summed over a 2D mesh the mesh points approach
$\bq = 0$ as $1/N$, and $\sum_\bq 1/q$ over an $N\times N$ mesh grows like $\ln N$: the
screening, and with it the gap, drifts logarithmically with the mesh instead of converging. His
numbers: the WSe$_2$ gap moved from 2.222 to 2.197~eV with the mesh; after projecting the charge
out of every fitted density, $d' = d - c\,(c^\dagger d)/(c^\dagger c)$, it is 2.233~eV and
stable; the 2D h-BN regression reference moved by $+0.92$~eV. The fix is in his PR.

The old large-cell scheme has the same fitted densities with the same spurious charge, but its
Coulomb matrix is not $2\pi/(Aq)$. $V(\bq)$ is a real-space lattice sum truncated at the
minimum-image radius $R_c \approx L/2$ of the supercell, whose charge-channel value is finite,
$\int_0^{R_c}\!2\pi r\,\mathrm{d}r/r = 2\pi R_c$ per cell area. The leak therefore does not
diverge; it produces a finite spurious screening proportional to $|m|^2 R_c/A \propto 1/L$,
i.e.\ a uniform downward shift of $\Sigma^c$ that decays only as the inverse cell size. That is
exactly the signature in the table: a gap that is too small, converges slowly from below, and
is still moving at 30~\AA. The sign of Wilhelm's fix ($+0.92$~eV for h-BN with his settings) and
the sign of our discrepancy (old scheme below) agree.

\section{Why the $\Gamma$-only scheme is immune}

The $\Gamma$-only implementation never lets a fitted density interact through the charge
channel of a sampled $V(\bq)$:
\begin{itemize}
\item the body of the dielectric matrix is built with the neutral projector,
      $V_b = \Pi V(0)\Pi$, $\Pi = 1 - \hat c\hat c^{\mathrm T}$, so the charge direction of the
      auxiliary basis carries no Coulomb interaction at all in the body;
\item the head and wings of $\varepsilon$ in the charge channel are built from the exact
      velocity-gauge transition dipoles $\langle a|[H,\br]|i\rangle/(\epsilon_a-\epsilon_i)$,
      checked by the Thomas--Reiche--Kuhn sum rule, and integrated analytically over the whole
      small Brillouin zone with the $2\pi/q$ kernel;
\item the regular $W(\Gamma)$ uses $V(0)$ with only the divergent $G = 0$ constant removed, and
      the charge direction of that matrix is quenched in $\varepsilon^{-1}$ (its Cholesky
      factorisation drops the negative mode).
\end{itemize}
This is the same treatment of the charge channel as Wilhelm's $\Gamma$ model, which is why the
two codes agree once his leak is removed. The agreement (15~meV at $8\times8$ versus
$24\times24$, with the remaining difference of the size expected from the $8\times8$ cell) is
the actual validation of the $\Gamma$-only scheme; the old large-cell numbers are not.

\section{Consequences}

\begin{enumerate}
\item Band gaps of 2D systems produced by \texttt{tensor\_large\_cell\_Gamma} or
      \texttt{ri\_rs\_large\_cell\_Gamma} with \texttt{KPOINTS\_W} $n\,n\,1$ ($n>1$) carry a
      cell-size-dependent error of order $0.1$~eV at 30~\AA\ and more than 1~eV on cells of
      20~\AA, and should not be used as references or regression values.
\item The $0.12$~eV difference between the two large-cell schemes at $14\times14$ is the old
      scheme's residual leak, consistent with the independent full-k-point result.
\item For 2D the large-cell code should be run $\Gamma$-only (\texttt{KPOINTS\_W 1 1 1}), which
      activates the head-and-wing correction; the same now holds for 3D cells
      (\texttt{PERIODIC XYZ}), where $V(0)$ is an Ewald sum with the $G=0$ term omitted.
\end{enumerate}

\paragraph{Data.} All inputs and outputs are in \texttt{large\_cell\_gamma/tests/881} of the
branch \texttt{rirs\_periodic\_gamma} (\texttt{github.com/RitajTyagi/cp2k}); the full-k
reference runs and the state-by-state table are in \texttt{881/kpoint\_ref\_pr6447/RESULTS.md}.
\end{document}
